% GENERATED FILE. Edit canonical lessons, then run npm run generate:books.
% canonical-source-hash: fnv1a64:c7dc786902ec7524
% canonical-lessons: PT-C452-barbeiro, PT-C452-aparar, PT-C452-esteticista, PT-C452-oculista, PT-C452-nodoa

\chapter{Cabelo e óculos}
\label{ch:pt-cabelo-e-oculos}
\hlchaptermodality{452}

\begin{chapteropening}
\textbf{By the end of this chapter:} \emph{I can use personal services: the barber, trimming, the beautician, the optician and a stain.}

Complete the last lesson of chapter 452: I can use personal services: the barber, trimming, the beautician, the optician and a stain.
\end{chapteropening}

\section[o barbeiro]{o barbeiro — the barber}

\label{lesson:PT-C452-barbeiro}



\begin{quote}
Before the new one: say the Portuguese for the locksmith, then the Portuguese for to patch.
\end{quote}

\subsection*{You'll want to know: o barbeiro}
\textbf{o barbeiro} — \textquotedblleft{}the barber\textquotedblright{}.

\begin{quote}
O barbeiro corta o cabelo ao meu avô há trinta anos.
\end{quote}

\begin{quote}
\emph{The barber has cut my grandfather's hair for thirty years.}
\end{quote}

Built on \emph{a barba}.

\begin{cousinweb}
\emph{barbeiro} is built from \emph{barba} 'beard' and \emph{-eiro} (EtymDB); \emph{barba} is inherited from Latin \emph{barba}.
\end{cousinweb}

\begin{grammarlens}[title={What you've built}]
One more word for telling, explaining and describing at B1.
\end{grammarlens}

\subsection*{Your turn}
\begin{itemize}
\raggedright
  \item \emph{Say it:} \emph{o barbeiro}
  \item \emph{Say it:} \emph{o barbeiro}, once more
  \item \emph{Say it:} a sentence of your own with \emph{o barbeiro}
  \item \emph{Recall:} say the Portuguese for the locksmith, then the Portuguese for to patch, then say \emph{o barbeiro} again
\end{itemize}

\begin{tcolorbox}[breakable,colback=teal!4,colframe=teal!35!black,title={Before you move on}]
What does o barbeiro mean? (\textquotedblleft{}The barber\textquotedblright{}.) Say it once more.
\end{tcolorbox}
\section[aparar]{aparar — to trim; to catch (something falling)}

\label{lesson:PT-C452-aparar}



\begin{quote}
Before the new one: say the Portuguese for to patch, then the Portuguese for the barber.
\end{quote}

\subsection*{You'll want to know: aparar}
\textbf{aparar} — \textquotedblleft{}to trim; to catch (something falling)\textquotedblright{}.

\begin{quote}
Só quero aparar a barba.
\end{quote}

\begin{quote}
\emph{I only want my beard trimmed.}
\end{quote}

\emph{Aparar a bola} — to catch the ball.

\begin{cousinweb}
\emph{aparar} goes back to Latin \emph{apparāre} 'to make ready', built on \emph{parāre} 'to prepare' (Lewis \& Short).
\end{cousinweb}

\begin{grammarlens}[title={What you've built}]
One more word for telling, explaining and describing at B1.
\end{grammarlens}

\subsection*{Your turn}
\begin{itemize}
\raggedright
  \item \emph{Say it:} \emph{aparar}
  \item \emph{Say it:} \emph{aparar}, once more
  \item \emph{Say it:} a sentence of your own with \emph{aparar}
  \item \emph{Recall:} say the Portuguese for to patch, then the Portuguese for the barber, then say \emph{aparar} again
\end{itemize}

\begin{tcolorbox}[breakable,colback=teal!4,colframe=teal!35!black,title={Before you move on}]
What does aparar mean? (\textquotedblleft{}To trim; to catch (something falling)\textquotedblright{}.) Say it once more.
\end{tcolorbox}
\section[a esteticista]{a esteticista — the beautician}

\label{lesson:PT-C452-esteticista}



\begin{quote}
Before the new one: say the Portuguese for the barber, then the Portuguese for to trim.
\end{quote}

\subsection*{You'll want to know: a esteticista}
\textbf{a esteticista} — \textquotedblleft{}the beautician\textquotedblright{}.

\begin{quote}
A esteticista trabalha no salão ao lado.
\end{quote}

\begin{quote}
\emph{The beautician works in the salon next door.}
\end{quote}

Built on \emph{a estética}, beauty care.

\begin{cousinweb}
\emph{esteticista} is built on \emph{estética}, a learned borrowing from Greek \emph{aisthētikós} 'of the senses', on \emph{aisthánesthai} 'to perceive' (EtymDB).
\end{cousinweb}

\begin{grammarlens}[title={What you've built}]
One more word for telling, explaining and describing at B1.
\end{grammarlens}

\subsection*{Your turn}
\begin{itemize}
\raggedright
  \item \emph{Say it:} \emph{a esteticista}
  \item \emph{Say it:} \emph{a esteticista}, once more
  \item \emph{Say it:} a sentence of your own with \emph{a esteticista}
  \item \emph{Recall:} say the Portuguese for the barber, then the Portuguese for to trim, then say \emph{a esteticista} again
\end{itemize}

\begin{tcolorbox}[breakable,colback=teal!4,colframe=teal!35!black,title={Before you move on}]
What does a esteticista mean? (\textquotedblleft{}The beautician\textquotedblright{}.) Say it once more.
\end{tcolorbox}
\section[o oculista]{o oculista — the optician}

\label{lesson:PT-C452-oculista}



\begin{quote}
Before the new one: say the Portuguese for to trim, then the Portuguese for the beautician.
\end{quote}

\subsection*{You'll want to know: o oculista}
\textbf{o oculista} — \textquotedblleft{}the optician\textquotedblright{}.

\begin{quote}
O oculista disse que preciso de óculos novos.
\end{quote}

\begin{quote}
\emph{The optician said I need new glasses.}
\end{quote}

Also the shop: \emph{ir ao oculista}.

\begin{cousinweb}
\emph{oculista} is built on \emph{óculo} 'eyeglass', from Latin \emph{oculus} 'eye' (Lewis \& Short).
\end{cousinweb}

\begin{grammarlens}[title={What you've built}]
One more word for telling, explaining and describing at B1.
\end{grammarlens}

\subsection*{Your turn}
\begin{itemize}
\raggedright
  \item \emph{Say it:} \emph{o oculista}
  \item \emph{Say it:} \emph{o oculista}, once more
  \item \emph{Say it:} a sentence of your own with \emph{o oculista}
  \item \emph{Recall:} say the Portuguese for to trim, then the Portuguese for the beautician, then say \emph{o oculista} again
\end{itemize}

\begin{tcolorbox}[breakable,colback=teal!4,colframe=teal!35!black,title={Before you move on}]
What does o oculista mean? (\textquotedblleft{}The optician\textquotedblright{}.) Say it once more.
\end{tcolorbox}
\section[a nódoa]{a nódoa — the stain}

\label{lesson:PT-C452-nodoa}



\begin{quote}
Before the new one: say the Portuguese for the beautician, then the Portuguese for the optician.
\end{quote}

\subsection*{You'll want to know: a nódoa}
\textbf{a nódoa} — \textquotedblleft{}the stain\textquotedblright{}.

\begin{quote}
Tenho uma nódoa de vinho na camisa.
\end{quote}

\begin{quote}
\emph{I have a wine stain on my shirt.}
\end{quote}

In Brazil: \emph{a mancha}. \emph{Tirar uma nódoa} — to get a stain out.

\begin{cousinweb}
\emph{nódoa} is inherited from Latin \emph{notula} 'little mark' (EtymDB), a small form of \emph{nota} 'mark'.
\end{cousinweb}

\begin{grammarlens}[title={What you've built}]
One more word for telling, explaining and describing at B1.
\end{grammarlens}

\subsection*{Your turn}
\begin{itemize}
\raggedright
  \item \emph{Say it:} \emph{a nódoa}
  \item \emph{Say it:} \emph{a nódoa}, once more
  \item \emph{Say it:} a sentence of your own with \emph{a nódoa}
  \item \emph{Recall:} say the Portuguese for the beautician, then the Portuguese for the optician, then say \emph{a nódoa} again
\end{itemize}

\begin{tcolorbox}[breakable,colback=teal!4,colframe=teal!35!black,title={Before you move on}]
What does a nódoa mean? (\textquotedblleft{}The stain\textquotedblright{}.) Say it once more.
\end{tcolorbox}
